\section*{第 \ref{ch:g10:universal-dna} 章——DNA：一种通用的遗传分子}

\begin{solution}{exo:g10:universal-dna:1}
磷酸基团、糖（脱氧核糖）和碱基。碱基有：腺嘌呤、
胸腺嘧啶、鸟嘌呤、胞嘧啶。
\end{solution}

\begin{solution}{exo:g10:universal-dna:2}
\texttt{ATGCCTAAGT}，逐碱基上下对齐书写。
\end{solution}

\begin{solution}{exo:g10:universal-dna:3}
C $= 21\%$；A $=$ T $= (100 - 42)/2 = 29\%$，各占。
\end{solution}

\begin{solution}{exo:g10:universal-dna:4}
\omterm{def:g10:universal-dna:gene}{基因}是 \omterm{def:g10:universal-dna:information}{DNA} 上的一段序列，其碱基序列携带合成一种\omterm{def:g10:chemistry-of-life:families}{蛋白质}的指令，
位于\omterm{def:g10:universal-dna:information}{染色体}的固定位置。\omterm{def:g10:universal-dna:gene}{等位基因}是某一\omterm{def:g10:universal-dna:gene}{基因}的
不同版本之一，与其他版本相比，只有少数几个碱基的
序列不同。
\end{solution}

\begin{solution}{exo:g10:universal-dna:5}
一种生物，其 \omterm{def:g10:universal-dna:information}{DNA} 中被插入了来自另一物种的\omterm{def:g10:universal-dna:gene}{基因}，
且能够读取并传递该\omterm{def:g10:universal-dna:gene}{基因}。例如：携带水母 GFP \omterm{def:g10:universal-dna:gene}{基因}的小鼠；
携带人胰岛素\omterm{def:g10:universal-dna:gene}{基因}的 \emph{E.~coli}。
\end{solution}

\begin{solution}{exo:g10:universal-dna:6}
$1500 \times 0.34 = \qty{510}{nm}$；$1500/10 = 150$ 圈。
\end{solution}

\begin{solution}{exo:g10:universal-dna:7}
$2 \times \num{3.2e9} \times \qty{0.34}{nm} \approx \qty{2.2}{m}$，
约为\omterm{def:g10:cells-common-unit:organelle}{细胞核}直径的 $\num{3.6e5}$ 倍。
\end{solution}

\begin{solution}{exo:g10:universal-dna:8}
“纯化”的组分中可能仍含有其他分子的痕迹，
因此仅凭其活性并不能证明 \omterm{def:g10:universal-dna:information}{DNA} 起作用。
在完整提取物中专门破坏 \omterm{def:g10:universal-dna:information}{DNA}，若活性随之消失，
而分解\omterm{def:g10:chemistry-of-life:families}{蛋白质}和 RNA 的酶对其没有影响，
则可确定该效应只来自 \omterm{def:g10:universal-dna:information}{DNA}。
\end{solution}

\begin{solution}{exo:g10:universal-dna:9}
若序列是固定重复的，则每个物种的碱基组成都相同（各碱基占 25\%），
但夏加夫发现，不同物种之间 A+T 的比例差异很大：碱基
的排列并非固定模式，序列可以变化，
这正是它能够携带信息的原因。
\end{solution}

\begin{solution}{exo:g10:universal-dna:10}
序列按顺序被读取以合成\omterm{def:g10:chemistry-of-life:families}{蛋白质}；若决定某个\omterm{def:g10:chemistry-of-life:families}{氨基酸}的位置上的碱基发生改变，
则可能使一种\omterm{def:g10:chemistry-of-life:families}{氨基酸}被另一种\omterm{def:g10:chemistry-of-life:families}{氨基酸}替换，
从而改变\omterm{def:g10:chemistry-of-life:families}{蛋白质}的形状或活性。若改变的碱基
位于\omterm{def:g10:universal-dna:gene}{基因}中不被读取的部分，或决定的是同一种\omterm{def:g10:chemistry-of-life:families}{氨基酸}，
则\omterm{def:g10:chemistry-of-life:families}{蛋白质}不变。
\end{solution}

\begin{solution}{exo:g10:universal-dna:11}
小鼠的所有\omterm{def:g10:cells-common-unit:cell}{细胞}都是由受精卵经分裂产生的，
每次分裂时，\omterm{def:g10:universal-dna:information}{DNA}（包括转\omterm{def:g10:universal-dna:gene}{基因}）都会被复制
（\omterm{prop:g10:universal-dna:helix}{互补性}），并分配到子\omterm{def:g10:cells-common-unit:cell}{细胞}中。因此皮肤\omterm{def:g10:cells-common-unit:cell}{细胞}、肝\omterm{def:g10:cells-common-unit:cell}{细胞}和精子都携带该\omterm{def:g10:universal-dna:gene}{基因}，
精子还会将其传给下一代。
\end{solution}

\begin{solution}{exo:g10:universal-dna:12}
$\num{4.6e6} \times \qty{0.34}{nm} \approx \qty{1.6}{mm}$，约为\omterm{def:g10:cells-common-unit:cell}{细胞}长度的 800
倍。但该分子的宽度只有 \qty{2}{nm}，
其体积约为 $\pi \times 1^2 \times \num{1.6e6}\,
\unit{nm^3} \approx \num{5e-3}\,\unit{\micro m^3}$，约占\omterm{def:g10:cells-common-unit:cell}{细胞}体积的 1\%；
折叠盘绕后，它可以绰绰有余地容纳在\omterm{def:g10:cells-common-unit:cell}{细胞}内。
\end{solution}

\begin{solution}{exo:g10:universal-dna:13}
人类 \omterm{def:g10:universal-dna:information}{DNA} 的长度是其 700 倍，但\omterm{def:g10:universal-dna:gene}{基因}数量只多约 5 倍，
即大部分 \omterm{def:g10:universal-dna:information}{DNA} 并不是\omterm{def:g10:universal-dna:gene}{基因}。多出的 \omterm{def:g10:universal-dna:information}{DNA} 中，有一部分调控\omterm{def:g10:universal-dna:gene}{基因}何时被使用，
另一部分是功能未知的重复序列。碱基
数量衡量的是“文本”的长度，而非生物体的复杂程度，
\omterm{def:g10:universal-dna:gene}{基因}的数量，尤其是\omterm{def:g10:universal-dna:gene}{基因}之间如何组合与调控，
才更为重要。
\end{solution}

\begin{solution}{exo:g10:universal-dna:14}
这将与 \omterm{def:g10:universal-dna:information}{DNA} 的通用性相矛盾
（见\cref{prop:g10:universal-dna:universal}）。既然共用的分子和密码
都来自共同祖先，这样的物种就必须起源于
另一个独立的生命起源，其\omterm{def:g10:universal-dna:gene}{基因}无法被
我们的\omterm{def:g10:cells-common-unit:cell}{细胞}读取，两个品系之间的转\omterm{def:g10:universal-dna:gene}{基因}也
将无法实现。
\end{solution}

\begin{solution}{exo:g10:universal-dna:15}
细菌生产的产物与人的\omterm{def:g10:chemistry-of-life:families}{蛋白质}完全相同（猪胰岛素与人胰岛素相差一个\omterm{def:g10:chemistry-of-life:families}{氨基酸}，
且可能引发免疫反应），并且可以大量生产，
不含猪的病原体。监管机构仍会询问，
纯化后的\omterm{def:g10:chemistry-of-life:families}{蛋白质}是否不含细菌杂质，
以及是否正确折叠并具有活性。
\end{solution}

\begin{solution}{pb:g10:universal-dna:1}
\textbf{1.} $720 \times 0.34 \approx \qty{245}{nm}$；$720/10 = 72$
圈。

\textbf{2.} \texttt{TACTCATTTCCTCTTCTTG}。

\textbf{3.} T $= 27\%$；G $=$ C $= (100 - 54)/2 = 23\%$。

\textbf{4.} $720/238 \approx 3.0$：每个\omterm{def:g10:chemistry-of-life:families}{氨基酸}由三个碱基决定（
最后三个碱基标志着肽链的终止）。

\textbf{5.} 共 $4^{720}$ 种序列，约 $10^{433}$ 种。化学上允许
任意排列顺序，因此在所有这些可能性中
选定的特定顺序便是一条消息——即信息。

\textbf{6.} 通用性：所有生物都使用相同的四种\omterm{def:g10:universal-dna:nucleotide}{核苷酸}，
并按相同的规则读取序列。

\textbf{7.} 小鼠的所有\omterm{def:g10:cells-common-unit:cell}{细胞}都是由卵经分裂而来，
因此导入卵中的\omterm{def:g10:universal-dna:gene}{基因}会被复制到每一个\omterm{def:g10:cells-common-unit:cell}{细胞}中。
若导入成体皮肤，则它只会存在于
被处理的皮肤\omterm{def:g10:cells-common-unit:cell}{细胞}及其后代中。

\textbf{8.} 每种\omterm{def:g10:cells-common-unit:cell}{细胞}类型都含有完整的 \omterm{def:g10:universal-dna:information}{DNA}，包括转\omterm{def:g10:universal-dna:gene}{基因}；
\omterm{def:g10:cells-common-unit:cell}{细胞}分化并不会丢弃\omterm{def:g10:universal-dna:gene}{基因}。

\textbf{9.} 是的：生殖\omterm{def:g10:cells-common-unit:cell}{细胞}携带该\omterm{def:g10:universal-dna:gene}{基因}。若小鼠的一对\omterm{def:g10:universal-dna:information}{染色体}中
只有一条带有该\omterm{def:g10:universal-dna:gene}{基因}的一个拷贝，则约一半后代会继承它
并发光。

\textbf{10.} \omterm{def:g10:chemistry-of-life:families}{蛋白质}不是信息：它被用尽后不再补充，
因此荧光会逐渐消退。\omterm{def:g10:universal-dna:gene}{基因}在每次分裂时被复制，
并持续被读取，因此荧光是持久并可遗传的。

\textbf{11.} \omterm{def:g10:universal-dna:information}{DNA} 组分；用一种能破坏
\omterm{def:g10:universal-dna:information}{DNA} 的酶进行处理。

\textbf{12.} 细菌从另一菌株摄取 \omterm{def:g10:universal-dna:information}{DNA}，将其整合，
并表达和传递该性状：一个\omterm{def:g10:universal-dna:gene}{基因}从一个生物体转移到另一个生物体，
这正是\omterm{prop:g10:universal-dna:universal}{转基因技术}有意进行的操作。

\textbf{13.} A$\,=\,$T 与 G$\,=\,$C 支持配对；
A+T 的比例可变，说明碱基对的排列顺序是自由的。

\textbf{14.} A 和 G 为大的双环碱基，T 和 C 为小的单环碱基：
一大一小配对时距离总是相同，
因此两条骨架沿整个分子始终保持 \qty{2}{nm} 的间距。
若两个大碱基配对则会隆起，两个小碱基配对则会留下空隙。

\textbf{15.} 因为每条链都是另一条链的互补链：
将二者分开，再依配对规则在每条链上合成其互补链，
即得两个完全相同的分子。

\textbf{16.} $2 \times \num{2.7e9} \times \qty{0.34}{nm} \approx
\qty{1.8}{m}$.

\textbf{17.} $\qty{1.8}{m} / \qty{6}{\micro m} = \num{3e5}$：一个
三十万倍的压缩。

\textbf{18.} $720 / \num{5.4e9} \approx \num{1.3e-7}$：约
\omterm{def:g10:cells-common-unit:cell}{细胞} \omterm{def:g10:universal-dna:information}{DNA} 的千万分之一。

\textbf{19.} $\num{2e4} \times \num{3e4} = \num{6e8}$ 个碱基对，约
占\omterm{def:g10:universal-dna:gene}{基因}组的 22\%；其余部分位于\omterm{def:g10:universal-dna:gene}{基因}之间。

\textbf{20.} 小鼠\omterm{def:g10:cells-common-unit:cell}{细胞}约携带 \qty{1.8}{m} 的 \omterm{def:g10:universal-dna:information}{DNA}；由于
该 \omterm{def:g10:universal-dna:information}{DNA} 的构建与读取方式与水母的完全相同，因此
两万个\omterm{def:g10:universal-dna:gene}{基因}中那一个外源\omterm{def:g10:universal-dna:gene}{基因}的读取方式与其他\omterm{def:g10:universal-dna:gene}{基因}一样。
\end{solution}
